\documentclass{article}

\usepackage{tabularx}
\usepackage{booktabs}
\usepackage{longtable}
\usepackage{array}
\usepackage{enumitem}
\usepackage[margin=1in]{geometry}
\usepackage[round]{natbib}
\usepackage{hyperref}
\hypersetup{colorlinks=true, linkcolor=red, citecolor=blue, urlcolor=blue}

\setlist{nosep, leftmargin=*}

\title{Interview Report\\\projname}

\author{\authname}

\date{}

\input{../../Comments.text}
\input{../../Common.text}

\newcommand{\intlink}[3]{\hyperref[int:#1-#2]{#3}}
\newcommand{\intref}[2]{\intlink{#1}{#2}{#1, #2}}

% Link into the published SRS PDF (same convention as the other documents)
\newcommand{\srsurl}{https://ayushpatel2026.github.io/Keystone/SRS/SRS.pdf}
\newcommand{\srsref}[1]{\href{\srsurl\#nameddest=#1}{\texttt{\detokenize{#1}}}}

% Empathy map: caption, label, says, thinks, does, feels
\newcommand{\empathymap}[6]{%
    \begin{longtable}{|p{0.465\textwidth}|p{0.465\textwidth}|}
    \caption{#1}\label{#2}\\
    \hline
    \textbf{Says} & \textbf{Thinks}\\ \hline
    \endfirsthead
    \hline
    \textbf{Says} & \textbf{Thinks}\\ \hline
    \endhead
    #3 & #4 \\ \hline
    \textbf{Does} & \textbf{Feels}\\ \hline
    #5 & #6 \\ \hline
    \end{longtable}}

\begin{document}

    \maketitle

    \begin{table}[hp]
        \caption{Revision History} \label{TblRevisionHistory}
        \begin{tabularx}{\textwidth}{llX}
\toprule
\textbf{Date} & \textbf{Developer(s)} & \textbf{Change}\\
\midrule
Nov 16 2026 & Ayush Patel & First version of interview report \\
\bottomrule
        \end{tabularx}
    \end{table}

    \newpage
    \tableofcontents
    \newpage

    \section{Purpose}

    This report is the fall custom document for \projname{}. It describes the
    interviews we held with stakeholders of the capstone project selection and
    approval process. It describes what we asked, who we asked, what we heard, how we analysed
    it, and what it led to. The Software Requirements Specification (SRS) already
    traces each interview statement to the stakeholder requirements it shaped
    (SRS Appendix A). This report records the evidence behind those requirements
    and the reasoning from interview to requirement. It also records the
    limits of that evidence, so that readers can judge how far to trust it.

    \section{Methodology}

    \subsection{Interview type}

    We ran semi-structured interviews. We prepared a guide of questions for each
    stakeholder type and used it as a starting point. We did not ask every
    interviewee every question. Instead we followed what each person told us and
    probed areas where they had experience or strong opinions. This suited our
    goal, which was to learn how each stakeholder experiences the current process
    and which problems matter most to them, not to collect comparable answers to a
    fixed questionnaire.

    \subsection{Procedure}

    \begin{itemize}
        \item We held six interviews between October 1 and October 30, 2026,
        covering seven people (Section~\ref{sec:subjects}). One interview was a
        joint session with two external supervisors.
        \item At least two members of the team attended every interview, and all
        four attended most of them. At least two members took notes in each
        interview.
        \item The guide opened with a short statement of who we are, what
        \projname{} is meant to do, and the purpose of the interview. The statement
        said that there were no right or wrong answers and that we wanted honest
        experiences. The guide moved from the interviewee's overall experience to
        specific parts of the process, and then asked what works well and what they
        would change.
        \item Interviews were not audio or video recorded, and we did not make
        transcripts.
        \item The guide was written in early October and then revised after
        feedback from our TA on October 13 (Section~\ref{sec:questions}). The
        interview with Dr.~Smith (October 1) took place before the guide was
        written, and the internal supervisor interview (October 7) took place
        before the TA's feedback. Both were therefore less structured than the
        later interviews.
        \item The internal supervisor interview was held in person, all others were on Microsoft Teams.
    \end{itemize}

    \subsection{Data preparation}

    Because we have notes and no transcripts, the notes are our primary data.
    After each interview we cleaned the notes. We fixed wording for clarity and
    removed interviewee names for students and supervisors. We separated what
    the interviewee said from our own comments, which are marked as
    \emph{interviewer notes} in Section~\ref{sec:raw}. We then
    tagged each statement with a code for the stakeholder group and a number or short name.
    The codes are CI (CS instructor), S (student), ES (external supervisor),
    IS (internal supervisor), SEI (SE instructor), and TI (TRON instructor). The SRS stakeholder
    requirements cite these tags, which is how the requirements trace back to specific
    statements. The notes record what we understood the interviewee to say. They are not
    verbatim quotes, and the statements in this report are paraphrases.

    \subsection{Analysis method: empathy map}

    We analysed the notes with an empathy map, a design thinking tool that
    sorts what is known about a group of people into four quadrants
    \citep{gray2010gamestorming, gibbons2018empathy}:

    \begin{itemize}
        \item \textbf{Says:} what they stated directly.
        \item \textbf{Thinks:} the beliefs and priorities behind what they said.
        \item \textbf{Does:} what they reported doing in their own work.
        \item \textbf{Feels:} the reactions and concerns they expressed.
    \end{itemize}

    We chose an empathy map for three reasons. First, it separates what people
    said and did, which we observed, from what they think and feel, which we
    infer. That distinction matters for note-based data, because it keeps our
    interpretation visible and checkable. Second, it gives each stakeholder
    group the same structure, which makes it easy to compare groups whose
    priorities conflict. Third, it focuses on needs and behaviour rather than
    on solutions. Stakeholders often proposed solutions (a newsletter, a
    template), and the map helps us ask what need is behind each one.

    We built one map per stakeholder group (Section~\ref{sec:maps}) from the cleaned
    notes after the interviews. The Says and Does quadrants come directly from
    the notes. The Thinks and Feels quadrants are our interpretation of the
    notes, and we have only included an entry where the notes give a clear basis
    for it. We then grouped the content of all six maps into cross-cutting themes
    (Section~\ref{sec:themes}), looking for problems raised independently by
    several groups and for places where groups disagree.

    \section{Subjects}
    \label{sec:subjects}

    We interviewed one or two people from each stakeholder group in Table~\ref{tab:subjects}.
    Participants were not sampled randomly. We interviewed stakeholders who were
    available to us. Names of students and supervisors are withheld. The course
    instructors are named because the course they run identifies them.

    \begin{table}[h!]
        \centering
        \renewcommand{\arraystretch}{1.2}
        \begin{tabularx}{\textwidth}{l l >{\raggedright\arraybackslash}X l}
            \toprule
            \textbf{Tag} & \textbf{Date (2026)} & \textbf{Interviewee} & \textbf{ID}\\
            \midrule
            SEI & October 1 & Dr.~Spencer Smith, instructor of the Software Engineering (SE) capstone course and our supervisor & \intlink{SEI}{1}{SEI, 1}\\
            IS & October 7 & Internal supervisor: a McMaster professor who supervises teams in all three capstone courses and submits proposals each year. Interviewed in person. & \intlink{IS}{1}{IS, 1}\\
            S & October 13 & Student who chose a project from the supervisor list and went through the SE (GitLab merge request) approval process & \intlink{S}{1}{S, 1}\\
            ES & October 13 & Two external supervisors from an industry consulting company, interviewed together. It was their first time supervising capstone teams, and one had previously been a capstone judge. They supervised eight teams. & \intlink{ES}{1}{ES, 1}\\
            TI & October 27 & Dr.~Alan Wassyng, instructor of the Mechatronics (TRON) capstone course & \intlink{TI}{1}{TI, 1}\\
            CI & October 30 & Dr.~Mehdi Moradi, instructor of the Computer Science (CS) capstone course & \intlink{CI}{1}{CI, 1}\\
            \bottomrule
        \end{tabularx}
        \caption{Interview subjects}
        \label{tab:subjects}
    \end{table}

    We did not interview a TA, although the question guide has a TA section.
    TA input to this report comes only from our own TA's feedback on the
    question guide. This gap is discussed in Section~\ref{sec:limits}.

    \section{Questions}
    \label{sec:questions}

    \subsection{Questions asked in the session}

    We prepared a guide with a set of general questions and extra questions for each
    stakeholder type. The guide is recorded in the following Github issue:
    \href{https://github.com/Ayushpatel2026/Keystone/issues/3}{Draft Questions for Stakeholder Interviews}.
    The questions below are the guide as drafted. Our TA gave feedback on October 13, and we
    reworded some questions in response. The reworded versions were not saved
    separately, so this section shows the draft. The TA's feedback was to:

    \begin{itemize}
        \item start with open questions about the interviewee's experience, and ask
        about problems afterwards;
        \item avoid leading wording (for example, ask ``Was anything confusing?''
        instead of ``What was confusing?'') and balance questions about
        problems with questions about what went smoothly;
        \item order questions to match the interviewee's experience, and tell
        the interviewee what the next set of questions is about;
        \item ask supervisors about their experience first (for example, how many
        times they have supervised), and ask a follow-up about how they handled
        a misunderstanding;
        \item add questions about how supervisors are found and how TAs choose
        their teams, and shorten the instructor guide to fit a 20--30 minute
        interview.
    \end{itemize}

    \subsubsection*{Introduction (all interviewees)}

    The guide opened with a statement along these lines. We are a capstone team
    at McMaster developing a centralized platform to streamline project
    selection and approval, which currently uses several disconnected tools. The
    interview aims to understand the interviewee's experience, the challenges they
    have met, and their needs. There are no right or wrong answers.

    \subsubsection*{General questions}
    \begin{itemize}
        \item Can you describe your overall experience with the current capstone project selection and management process?
        \item Have you encountered situations where a lack of information, communication, or transparency caused difficulties?
        \item What information do you find yourself repeatedly searching for or requesting from others?
        \item If you could make one improvement to the current process, what would it be?
        \item Is there anything about the process that works particularly well and that you would not want changed?
    \end{itemize}

    \subsubsection*{Students}
    \textit{Project selection}
    \begin{itemize}
        \item Can you walk us through how you found your capstone project, from initially looking for opportunities to getting your project approved?
        \item How did you find out which projects were available, and how did you determine which ones interested you?
        \item Was the information in project descriptions sufficient to decide whether a project was a good fit? What was missing?
        \item What information about a project would be most useful when deciding whether to pursue it?
        \item Did you ever pursue a project and later find out it was taken, unavailable, or different from its description? What happened?
        \item How did you contact supervisors, and how long did they take to reply?
        \item What would make you feel more confident that a project and its supervisor are a good fit for your team?
    \end{itemize}
    \textit{Team selection}
    \begin{itemize}
        \item How did you find your teammates? Did you already know the people you wanted to work with, or did you have to find them?
        \item What factors matter most to you when choosing teammates?
        \item Did anyone you know end up in a mismatched team? What went wrong?
    \end{itemize}
    \textit{Project approval}
    \begin{itemize}
        \item (SE) What was the GitLab merge request process like? (Tron/CS) What was your approval process like? Was anything confusing?
        \item After approval, did you change anything (supervisor contact, team name, GitHub link)? How did you report it, and was it clear how?
    \end{itemize}

    \subsubsection*{Supervisors}
    \textit{Submitting a project}
    \begin{itemize}
        \item Can you walk us through how you currently propose a capstone project and get it approved?
        \item What aspects of submitting and managing a capstone project require the most time or administrative effort?
        \item What would help you describe a project accurately (templates, examples, guidance on scope and skills)?
        \item Have confidentiality, intellectual property, or feasibility requirements created challenges during approval?
        \item If your project was not chosen, how did you find out? What would you want to be told, and when?
        \item What would make you more or less likely to offer a project again next year?
    \end{itemize}
    \textit{Working with student teams}
    \begin{itemize}
        \item How do you communicate with students who express interest in your project?
        \item How do you decide which team or teams to work with when several are interested?
        \item What information do you need from students before agreeing to supervise a team?
        \item How do you stay informed about the status and progress of your project and team during the year?
        \item Have you experienced difficulties with the submission process, such as unclear requirements, missing information, or repeated communication with course instructors?
        \item Have students misunderstood the scope, requirements, or deliverables of your project? What happened?
        \item How easy is it to keep track of interested teams, communicate with them, and know whether a project has been taken?
        \item What communication or administrative tasks do you find repetitive?
    \end{itemize}

    \subsubsection*{Course instructors}
    \begin{itemize}
        \item Walk us through how you handle project selection and approval for your program, from the project list to confirmed teams.
        \item Which tools and records do you use to track projects and teams? Who has access?
        \item How do you know which projects your students have taken?
        \item What are the rules for your students (hardware requirement, working with other programs, team size)? Who enforces them?
        \item How do you coordinate with the other instructors? What do you need from them, and what do they need from you?
        \item How do you handle projects chosen by teams from more than one program?
        \item How do you communicate with supervisors during the year? What takes the most time?
        \item What criteria do you use to judge the feasibility, suitability, and academic value of a proposed project?
        \item Which parts of approval need your direct judgment, and which follow standard rules?
        \item How do you handle changes to project descriptions, supervisors, team members, or other details after approval?
        \item What information or reports would help you make administrative decisions or find issues earlier?
        \item Which parts of the current workflow would you want to keep, even with a centralized platform?
    \end{itemize}

    \subsubsection*{TAs}
    The guide asked TAs what capstone tasks they do, how long these take and which
    are repetitive, and what the most common student and supervisor mistakes or
    questions are. We did not hold a TA interview.

    \subsection{Coverage of topics}

    Because the interviews were semi-structured, each one covered a different
    subset of the guide. Table~\ref{tab:coverage} shows which topics each interview reached,
    based on the notes. A dash means the topic does not appear in the notes. It does not
    necessarily mean the topic was not raised.

    \begin{table}[h!]
        \centering
        \small
        \renewcommand{\arraystretch}{1.15}
        \begin{tabularx}{\textwidth}{>{\raggedright\arraybackslash}X c c c c c c}
            \toprule
            \textbf{Topic} & \textbf{SEI} & \textbf{IS} & \textbf{S} & \textbf{ES} & \textbf{TI} & \textbf{CI}\\
            \midrule
            Overall experience of the process & Yes & Yes & Yes & Yes & Yes & Yes\\
            Proposals and project descriptions & Yes & Yes & Yes & Yes & Part & Part\\
            Project availability and status & Part & -- & Yes & Yes & -- & Yes\\
            Team formation and student backgrounds & Part & Part & Yes & Part & Yes & Yes\\
            Supervisor--team communication and agreement & Yes & Part & Yes & Yes & -- & Part\\
            Approval process & Yes & Yes & Yes & -- & Yes & Yes\\
            Changes and records after approval & Yes & -- & -- & -- & -- & --\\
            Course information and timelines for supervisors & Yes & Yes & Part & Yes & Part & --\\
            Tools and records in use & Yes & -- & Part & Part & Yes & Yes\\
            Platform scope and what to keep & Yes & Yes & -- & Part & Yes & Yes\\
            \bottomrule
        \end{tabularx}
        \caption{Topics covered by each interview (Yes: covered in detail; Part: touched on; --: not in the notes)}
        \label{tab:coverage}
    \end{table}

    \section{Raw Data}
    \label{sec:raw}

    We have notes and no transcripts. This section gives the cleaned notes for each interview.
    The SRS (Appendix A) holds the same statements, tagged and numbered, and the stakeholder
    requirements cite them. The notes here keep more context (for example the interviewee's
    background and the sequence of the process they described). Text marked
    \emph{interviewer note} is our comment and not the interviewee's statement.

    \subsection{SE course instructor (Dr.~Smith), October 1}
    \label{int:SEI-1}

    \textbf{Current process}
    \begin{itemize}
        \item In early summer, Dr.~Smith emails industry contacts and faculty to ask for
        project ideas. Interested supervisors fill out an MS Form (title, organization, contact, 100--200 word
        description, confidentiality, IP concerns, whether the project involves hardware).
        \item The responses go to a spreadsheet, and a script turns them into a LaTeX list.
        Projects with hardware are labelled TRON, and the rest are labelled SE/CS. The resulting
        PDF is sent to SE, CS, and TRON students. A second email tells supervisors that
        students will contact them.
        \item Student steps are informal. Students form a team, then either pick a project
        and email its supervisor, or propose their own. Picking from the list means
        contacting and meeting the supervisor, who must agree to work with the team. Teams can be rejected and
        have to look for another project. Student-proposed projects need enough complexity and
        design for four or five final-year students, and real stakeholders. A supervisor is
        recommended when domain expertise is needed, and some teams open a request without one.
        \item Teams open a GitLab merge request with two JSON files (project and team information) and a
        Markdown project description. Once approved, CI merges the files and updates
        \texttt{TeamComposition.csv}, which is his one-stop record of every team's project,
        contacts, and links.
        \item He assigns TAs to teams with the TAs, in a JSON file that CI uses to update the team files.
    \end{itemize}

    \textbf{Pain points}
    \begin{itemize}
        \item He checks many items by hand and gives the same feedback to many teams: a
        non-generic team name, a public GitHub repo from the template with him added as collaborator
        and Pages enabled, appropriate custom documents, valid stakeholders, and a supervisor where one is needed.
        The GitLab checklist is not fully accurate, and teams do not use it. He wants to keep judging
        feasibility, complexity, and scope himself, based on experience.
        \item Teams change titles, repository names, supervisors, or descriptions after approval without opening
        another merge request, so \texttt{TeamComposition.csv} drifts out of date. Teams update only
        their custom document choice, because it is marked.
        \item He knows which projects SE students chose, but not which projects TRON and CS students
        chose, so he cannot send thank-you emails to supervisors whose projects were not selected.
        \item He sends four emails to supervisors: a call for projects, an initial contact, a request to confirm
        willingness to supervise (because some students claim an agreement that the supervisor does not
        recall), and a welcome with course information. The last one addresses many supervisor questions,
        but it arrives after teams and projects are settled.
        \item Descriptions vary in quality. Standards could help, but a burdensome form could scare away good
        project ideas.
        \item \emph{Interviewer note:} our own team's first project was described as a new application but was
        about 80\% built.
    \end{itemize}

    \textbf{Preferences for the new platform}
    \begin{itemize}
        \item Features should be modular and extensible, so that an instructor can turn them off and the platform
        can serve other courses and universities.
        \item The platform should keep the GitLab merge request workflow and integrate with Git for the SE course.
        \item AI agents should not make approvals. Summaries or automated replies to repeated issues
        could be considered. Algorithmic matching of teams to projects could be interesting if it can be turned
        off and the overhead is worth it.
    \end{itemize}

    \subsection{Internal supervisor, October 7}
    \label{int:IS-1}

    \begin{itemize}
        \item Mostly acts as a supervisor. He submits proposals each year, waits to see what else is on the list,
        and interviews and supervises teams in all three courses. He attends lectures when possible and helps
        with presentations and Expo posters. He often feels outside the course.
        \item Once a supervisor submits the proposal form, the timing of approval is not known to supervisors.
        There is no standard procedure: when a student should start thinking about capstone, how to
        find a team, and how to find a supervisor for their own idea.
        \item Project schedules and due dates are confusing. Teams often lack the time to finish a project and do
        proper user testing and validation. Supervisors with teams in several courses do not know that
        each course has different timelines and requirements. The three courses are not well coordinated.
        \item He is unsure how student-proposed projects will be handled. He would like to see them.
        \item Descriptions are 100--200 words. He finds standardizing them hard and is unsure what the standard
        should be. One team delivered something related to, but very different from, the proposal.
        \item He wants a low-effort way to submit feedback at the end. One instructor sends feedback emails,
        but he does not recall other instructors asking.
        \item He wants to know who the TAs are for his teams and wants student skills to guide discussions.
        He would like the capstone process in one place, but is unsure that works across professors.
        \item \emph{Interviewer note:} Dr.~Smith does send a schedule and asks for feedback several times,
        and we do not know whether the other two instructors do.
    \end{itemize}

    \subsection{Student, October 13}
    \label{int:S-1}

    The student chose a project from the supervisor list and went through the SE (GitLab) process.
    \begin{itemize}
        \item Formed a team of classmates over the summer and had project ideas before term. In hindsight, picking
        the project before the teammates might have given a better experience, since teammates with shared
        interests would be working on their preferred project.
        \item Each member chose their top five projects after reading them all, and the team settled on a top three
        in case some were unavailable.
        \item Project descriptions were often vague, and some were ignored and only the title was read.
        Information on hardware, software stack, and industry would help.
        \item One supervisor held a meeting with all interested teams right before the deadline, which added
        uncertainty. Another team competed for the same project, and losing it was frustrating. There was
        considerable uncertainty and worry about not having a project.
        \item A template, guidelines, or an FAQ for proposing their own project would reduce uncertainty about size,
        scope, and content requirements.
        \item They want to know when a project is taken, open, or dropped, and suggested a newsletter.
        \item Communication with supervisors is unclear (an agenda, who leads, how often). Meetings were useful for
        clearing a backlog of questions with an expert. Email is preferred for simple questions, and suggested
        meetings aligned with deliverables would help teams know whether they are on track.
        \item The GitLab process was intuitive because of prior GitHub/GitLab experience, and the JSON files
        were a hassle but not a problem.
        \item Finding where things are is time consuming and needs many open tabs. Google Calendar was helpful
        but would benefit from more information. Having everything in one place, like Avenue, would help.
    \end{itemize}

    \subsection{External supervisors, October 13}
    \label{int:ES-1}

    Two supervisors from an industry consulting company, interviewed together. It was their first time supervising
    capstone teams. One had been a capstone judge, so they knew a little about the course. They supervised eight
    teams.
    \begin{itemize}
        \item The proposal and the first meetings with a professor and the team went smoothly. Knowing project
        statuses was a challenge, and they wanted to know who is talking to whom.
        \item They wanted to keep selection fair. Each meeting with a team shared new information, so one of them
        sent summaries of the meetings to all the teams by hand. They suggested a shared place for meeting notes
        and a template for summaries.
        \item They were unsure when to reply to students, whether to rush, and whether to speak to everyone who
        emailed. With eight teams there was a lot of administration through emails and team meetings.
        Groups usually had one representative email with the members in CC, which made emails hard to search.
        With many names and no group numbers, they defaulted to the first person who emailed.
        \item A template with more structure would be better than an open-ended description box. It is hard to
        know what a capstone team is expected to do, so they took a best-effort approach. They wanted
        examples of recent capstone projects.
        \item They wanted the documents from the welcome email (course outline, course and project expectations,
        timelines, and deadlines) before submitting a proposal. More information up front would avoid
        back-and-forth with the instructor over processes, expectations, and IP.
        \item They wanted a schedule, and suggested meeting times, to clear up expectations. Some teams come early
        and some late.
        \item They knew nothing about students before the first meeting and relied on gut feeling. They looked up some
        students on LinkedIn, but some had no useful profile. They were confused about the difference between
        CS and SE students, and wanted to see a team bio and skills. They did not know that the final decision
        was theirs. They also wanted teams matched to projects by preferences and skills.
        \item After a supervisor meeting, students can also say no. The platform should record whether both
        sides agreed before the process goes to instructor approval.
        \item Being notified of tasks is fine, but having to remember to open the platform would be terrible.
        Email notifications and a checklist are preferred. Too many platforms to log in to is tedious, but a
        new application is fine.
        \item \emph{Interviewer note:} when students contact supervisors through the platform, discussion should
        stay in the context of an application or student post, like comments and replies on a pull request, and
        not full messaging. Meetings will happen outside the platform.
    \end{itemize}

    \subsection{TRON course instructor (Dr.~Wassyng), October 27}
    \label{int:TI-1}

    \begin{itemize}
        \item The TRON course was separated from the SE course because of growing class size. The TRON and CS
        courses run differently, and the TRON course is more flexible than SE. He emphasizes personal
        responsibility and meets each team informally every week. Teams meet in person during class.
        \item Most of the course runs on Microsoft Teams, with a private channel per team plus general channels
        and shared files. Avenue holds grades, submissions, and deliverables. Deliverables are due at the end
        of the week and accepted without penalty until the following Tuesday.
        \item The beginning of the capstone courses is the most similar part across courses, and they have a similar
        number of deliverables.
        \item Projects are approved on Teams: a project is approved when its team receives a group number.
        \item He receives the list of students, marks off those who are in a group by hand, and contacts the rest.
        He tries not to force students into groups, so some are delayed. One student joined very late and the
        situation went to the undergraduate advisor.
        \item Teams channels must be created and populated from Avenue or by hand, and students are removed if they
        are not on Avenue. Populating a platform with student information is not trivial. An Avenue course is
        optional and created on request.
        \item SE students on TRON projects are treated as TRON students, and their grades are sent to the SE
        instructor. TRON students cannot join software-only projects. ECE students on TRON projects must be graded
        by an ECE professor, and cross-discipline capstones such as iBioMed need extra involvement. Accreditation
        is a concern.
        \item Projects are collected over the summer through the SE instructor and his contacts, and kept in a
        LaTeX document. Many TRON projects are student-proposed, and these are often high quality and
        people-centred.
        \item The capstone is student-centred. Students own the project, even if stakeholders are involved, and
        the main rule is that students must be able to talk about it (presentations, Expo). Supervisor-only
        deliverables are secondary and should not overwhelm students. He sends supervisors a list of the previous
        year's deliverables.
        \item Open question: how involved the platform will be in the capstone process.
    \end{itemize}

    \subsection{CS course instructor (Dr.~Moradi), October 30}
    \label{int:CI-1}

    \begin{itemize}
        \item Projects from other professors are collected through the form that the SE instructor sends out.
        Most communication about projects happens in August. He would like to widen who receives the form
        for external projects.
        \item He wants project state kept up to date automatically, since supervisors currently have to report it.
        A project accepted by a team should become unavailable.
        \item Listed projects are accepted automatically since he has already seen them. More than half of teams
        define their own project, which is handled on a discussion board. A board of proposals that he can approve
        would be less disorganized.
        \item Students struggle with team formation, which relies on talking to friends. Features that help students
        communicate would help.
        \item He does not use GitLab. Students keep GitHub repositories and share them with him. Assignment submission
        should stay where it is, and instructor and TA access should be configurable.
        \item He does not want engagement past the selection stage and expects supervisors to contact him if there
        is a problem. Features should be optional, intuitive, and not create extra work.
        \item Open question: whether data can stay internal to the university when external supervisors are involved,
        and whether internal hosting would prevent them from using the platform.
        \item He also asked for a way to book two meetings per term with teams and remind teams of missed meetings,
        supervisor--instructor communication for each project, and record keeping between students and supervisors
        or instructors.
    \end{itemize}

    \section{Analysis}

    \subsection{Empathy maps}
    \label{sec:maps}

    The maps below summarize each stakeholder group. Says and Does come directly from the notes in
    Section~\ref{sec:raw}. Thinks and Feels are our interpretation of those notes.

    \empathymap{Empathy map: student (\intref{S}{1})}{tab:map-student}
    {\begin{itemize}
         \item Descriptions are vague, and some projects were judged by title only.
         \item Wants to know a project's hardware, software stack, and industry.
         \item Wants to know when a project is taken, open, or dropped.
         \item Wants a template and guidelines for proposing their own project.
         \item Communication with supervisors is unclear. Suggested meetings aligned to deliverables would help.
         \item Wants everything in one place.
    \end{itemize}}
    {\begin{itemize}
         \item Picking a project before choosing teammates might have worked better.
         \item A ranked backup list is needed because projects may be unavailable.
         \item Meetings with an expert are valuable for closing knowledge gaps.
         \item The GitLab approval process was intuitive, and the JSON files were a minor hassle.
    \end{itemize}}
    {\begin{itemize}
         \item Formed a team with known classmates over the summer.
         \item Each member read all projects and chose a top five, and the team agreed a top three.
         \item Competed with another team for a project.
         \item Used many tabs and Google Calendar to find information.
    \end{itemize}}
    {\begin{itemize}
         \item Uncertain and worried about ending up without a project.
         \item Frustrated at losing a project to another team.
         \item Unsure whether supervisors would have time to meet.
         \item Unsettled by a supervisor's meeting with all interested teams just before the deadline.
    \end{itemize}}

    \empathymap{Empathy map: external supervisors (\intref{ES}{1})}{tab:map-es}
    {\begin{itemize}
         \item Wants a template rather than an open box, examples of past projects, and course information
         before proposing.
         \item Wants to know student backgrounds, a team bio, and who is talking to whom.
         \item Wants notifications and a checklist so they do not have to remember to check a platform.
         \item Wants the agreement of both sides recorded before instructor approval.
         \item Wants a shared place for meeting notes and project status.
    \end{itemize}}
    {\begin{itemize}
         \item Selection should be fair, and information shared with one team should reach all teams.
         \item Information about the course would reduce back-and-forth with the instructor.
         \item A new platform is acceptable, but another place to log in to is a cost.
         \item They believed they had to accept whichever team asked, until told otherwise.
    \end{itemize}}
    {\begin{itemize}
         \item Supervised eight teams through email and team meetings.
         \item Sent summaries of team meetings to all teams by hand.
         \item Looked up students on LinkedIn.
         \item Replied to the first person who emailed when groups had many names and no group number.
         \item Wrote proposals on a best-effort basis.
    \end{itemize}}
    {\begin{itemize}
         \item Lost about the timeline: how fast to reply and whether to talk to everyone.
         \item Took a leap of faith in agreeing to a team they knew nothing about.
         \item Burdened by the administration of managing eight teams.
         \item Confused about the difference between CS and SE students.
    \end{itemize}}

    \empathymap{Empathy map: internal supervisor (\intref{IS}{1})}{tab:map-is}
    {\begin{itemize}
         \item The three courses are not well coordinated, and their timelines differ.
         \item The timing and method of approval are not known to supervisors.
         \item There is no standard procedure, for example for finding a supervisor for a student idea.
         \item Wants a low-effort way to give end-of-term feedback and wants to know each team's TA.
         \item Wants the capstone process in one place, if that works with other professors.
    \end{itemize}}
    {\begin{itemize}
         \item Standardizing project descriptions is hard, and it is unclear what the standard should be.
         \item Teams often lack the time to finish and to validate properly.
         \item Student skills would help guide discussions.
         \item The supervisor's role is on the outside of the course.
    \end{itemize}}
    {\begin{itemize}
         \item Submits proposals each year, then waits to see what else is listed.
         \item Interviews and supervises teams in all three courses.
         \item Attends lectures, presentations, and Expo.
         \item Sends feedback on each team to the instructors without being asked.
    \end{itemize}}
    {\begin{itemize}
         \item Confused about what is due and when.
         \item Unsure how student-proposed projects will be handled.
         \item Not sure that other professors would agree to a single platform.
    \end{itemize}}

    \empathymap{Empathy map: SE course instructor (\intref{SEI}{1})}{tab:map-sei}
    {\begin{itemize}
         \item The same feedback goes to many teams.
         \item Wants to know which projects were not selected, in order to thank their supervisors.
         \item Wants coordination and records generated automatically.
         \item Wants the platform to be modular and to keep the GitLab merge request workflow.
         \item No AI agents making approvals.
    \end{itemize}}
    {\begin{itemize}
         \item Judging feasibility, scope, and complexity needs his experience and should stay with him.
         \item Pull requests teach students about CI and allow discussion with students.
         \item Standards could discourage good projects, so optional questions may help.
         \item Matching is interesting only if it can be turned off and the overhead is worth it.
    \end{itemize}}
    {\begin{itemize}
         \item Recruits projects each summer and runs the form, script, and four emails.
         \item Reviews each merge request by hand and assigns TAs with the TAs.
         \item Uses \texttt{TeamComposition.csv} as his single record.
         \item Confirms with supervisors that they agreed to supervise a team.
    \end{itemize}}
    {\begin{itemize}
         \item Frustrated at spending time on repeated feedback.
         \item Concerned that records drift out of date after approval.
         \item Concerned that supervisors get key course information too late.
         \item Wary of putting burden on supervisors with a long form.
    \end{itemize}}

    \empathymap{Empathy map: TRON course instructor (\intref{TI}{1})}{tab:map-ti}
    {\begin{itemize}
         \item Many TRON projects are proposed by students, and they are often high quality.
         \item Populating a platform with student information is not trivial.
         \item Cross-discipline rules are complex: SE, ECE, iBioMed.
         \item Supervisor-only deliverables should not overwhelm students.
         \item Open question: how involved the platform will be.
    \end{itemize}}
    {\begin{itemize}
         \item The capstone belongs to the students, and students must be able to talk about their project.
         \item The course should stay flexible and run on personal responsibility.
         \item Accreditation matters in cross-discipline projects.
    \end{itemize}}
    {\begin{itemize}
         \item Runs the course on Teams, with Avenue for grades and submissions.
         \item Marks students off a roster by hand and contacts those without a group.
         \item Approves a project when a team receives a group number.
         \item Meets each team informally every week.
         \item Sends supervisors the previous year's deliverables.
    \end{itemize}}
    {\begin{itemize}
         \item Reluctant to force students into groups.
         \item Concerned about students who join groups late.
         \item Uncertain how much a platform would change a course that already works for him.
    \end{itemize}}

    \empathymap{Empathy map: CS course instructor (\intref{CI}{1})}{tab:map-ci}
    {\begin{itemize}
         \item Wants project state kept up to date automatically, and wants taken projects shown as unavailable.
         \item Wants a board of proposals he can approve in place of the discussion board.
         \item Students struggle with team formation.
         \item Features should be optional, intuitive, and not add work.
         \item Does not want engagement past the selection stage.
         \item Asked whether data can stay internal when external supervisors are involved.
    \end{itemize}}
    {\begin{itemize}
         \item Listed projects need no review because he has already seen them.
         \item Supervisors will contact him if a problem arises.
         \item Assignment submission should stay where it is.
    \end{itemize}}
    {\begin{itemize}
         \item Collects projects through the SE instructor's form, mostly in August.
         \item Handles the more than half of teams with their own project through a discussion board.
         \item Does not use GitLab. Students share GitHub repositories with him.
    \end{itemize}}
    {\begin{itemize}
         \item Unwilling to take on extra work from a new platform.
         \item Disorganized by the discussion board approach.
         \item Uncertain about whether hosting and external users can work together.
    \end{itemize}}

    \subsection{Themes}
    \label{sec:themes}

    Comparing the maps gave the ten themes below. For each, we list which interviews raised it.

    \begin{enumerate}
        \item \textbf{Project descriptions are inconsistent and do not set expectations.}
        Students ignore vague descriptions and want hardware, stack, and industry. Supervisors do not
        know what a proposal should contain, and one team's project differed from its description.
        Raised in \intref{S}{1}, \intref{ES}{1}, \intref{IS}{1}, \intref{SEI}{1}.
        \textit{Tension:} the SE instructor and internal supervisor both warn that
        strict standards could scare away good ideas.

        \item \textbf{Nobody can see which projects are open, taken, or unselected.}
        Students worry about taken projects, supervisors cannot tell project statuses, and the instructors cannot
        tell which projects other courses' students chose. Raised in \intref{S}{1}, \intref{ES}{1},
        \intref{CI}{1}, \intref{SEI}{1}.

        \item \textbf{Student-proposed projects have no clear path and are hard to see.}
        Many TRON and CS teams, more than half in CS, propose their own project, but these are
        discussed on a board that supervisors cannot see. Raised in \intref{IS}{1}, \intref{TI}{1},
        \intref{CI}{1}, \intref{S}{1}.

        \item \textbf{Teammates and students' backgrounds are unknown.}
        Students rely on friends to form teams, and supervisors know nothing about the students
        before meeting them. Raised in \intref{S}{1}, \intref{CI}{1}, \intref{ES}{1}, \intref{IS}{1},
        \intref{TI}{1}.

        \item \textbf{Supervisors get course guidance and schedules too late.}
        Supervisors want the course outline, expectations, timelines, and examples before proposing. They also
        do not know how timelines differ across courses. Raised in \intref{ES}{1}, \intref{IS}{1},
        \intref{SEI}{1}, \intref{TI}{1}.

        \item \textbf{Communication during selection is repetitive, and agreements are not recorded.}
        Supervisors get many emails, share meeting information by hand, and cannot tell who is talking to
        whom. Students say a supervisor agreed that the supervisor does not recall. Raised in
        \intref{ES}{1}, \intref{SEI}{1}, \intref{S}{1}.

        \item \textbf{Too many tools, and users do not want to remember to check another one.}
        Students and supervisors want one place, and notifications and a checklist in place of checking.
        Raised in \intref{S}{1}, \intref{ES}{1}, \intref{IS}{1}.

        \item \textbf{Instructor review is repetitive, and records drift after approval.}
        The SE instructor repeats the same feedback and records go stale. Instructors also populate rosters
        and track who is in a group by hand. Raised in \intref{SEI}{1}, \intref{TI}{1}, \intref{CI}{1}.

        \item \textbf{Courses differ, so the platform must be configurable.}
        SE approves by merge request, TRON by group number on Teams, and CS accepts listed projects
        automatically. Team rules and tools also differ. Raised in \intref{SEI}{1}, \intref{TI}{1},
        \intref{CI}{1}.

        \item \textbf{Automation and platform involvement must stay limited.}
        Instructors want to keep approval decisions and feasibility judgments. The CS instructor wants
        no extra work, the TRON instructor wants the course to remain student-centred, and supervisors
        already feel the administrative load. Raised in \intref{SEI}{1}, \intref{CI}{1}, \intref{TI}{1},
        \intref{ES}{1}.
    \end{enumerate}

    \subsection{How the themes were used}

    Each theme led to stakeholder requirements in the SRS, and these to functional requirements
    (Table~\ref{tab:trace}). The SRS stakeholder requirements cite the specific interview statements.

    \begin{longtable}{>{\raggedright\arraybackslash}p{0.06\textwidth} >{\raggedright\arraybackslash}p{0.43\textwidth} >{\raggedright\arraybackslash}p{0.43\textwidth}}
        \caption{From themes to SRS requirements}\label{tab:trace}\\
        \toprule
        \textbf{Theme} & \textbf{Stakeholder requirements} & \textbf{Functional requirements}\\
        \midrule
        \endfirsthead
        \toprule
        \textbf{Theme} & \textbf{Stakeholder requirements} & \textbf{Functional requirements}\\
        \midrule
        \endhead
        1 & \srsref{SR_ProposalTemplate}, \srsref{SR_OptionalFields}, \srsref{SR_ProposalGuidance} & \srsref{FR_SupervisorProposal}, \srsref{FR_ProposalOptionalFields}\\ \hline
        2 & \srsref{SR_BrowseProjects}, \srsref{SR_ProjectStatus}, \srsref{SR_StatusNotify}, \srsref{SR_UnselectedProjects} & \srsref{FR_ViewApprovedProjects}, \srsref{FR_ViewProjectStatus}, \srsref{FR_ReadCourseInformation}\\ \hline
        3 & \srsref{SR_StudentProposals} & \srsref{FR_StudentProposal}\\ \hline
        4 & \srsref{SR_StudentProfiles}, \srsref{SR_TeamContacts} & \srsref{FR_StudentProfile}, \srsref{FR_TeamProfile}, \srsref{FR_TeamDiscovery}, \srsref{FR_AssignTA}\\ \hline
        5 & \srsref{SR_CourseInfoEarly}, \srsref{SR_ProposalGuidance}, \srsref{SR_CrossCourseTimelines} & \srsref{FR_ReferenceMaterial}\\ \hline
        6 & \srsref{SR_SinglePlace}, \srsref{SR_TeamContacts}, \srsref{SR_ReviewDialogue} & \srsref{FR_ProjectInterest}, \srsref{FR_DeclineTeam}, \srsref{FR_ProjectThreads}\\ \hline
        7 & \srsref{SR_PushNotify}, \srsref{SR_Checklist}, \srsref{SR_SinglePlace} & \srsref{FR_ThreadsNotify}, \srsref{FR_RequiredActionsList}, \srsref{FR_ActionRequiredEmail}\\ \hline
        8 & \srsref{SR_TeamSubmissionChecks}, \srsref{SR_TeamValidity}, \srsref{SR_ReduceAdmin}, \srsref{SR_TAApproval}, \srsref{SR_Enrollment} & \srsref{FR_ApprovalPreflightValidate}, \srsref{FR_ApprovalRevision}, \srsref{FR_ExportCourseInformation}, \srsref{FR_ImportEnrollment}, \srsref{FR_NeedsGroup}\\ \hline
        9 & \srsref{SR_CourseConfig}, \srsref{SR_Enrollment} & \srsref{FR_DisableFeatures}, \srsref{FR_ExternalApproval}, \srsref{FR_ValidateTeamRules}\\ \hline
        10 & \srsref{SR_HumanApproval}, \srsref{SR_OptionalFields}, \srsref{SR_ReviewDialogue} & \srsref{FR_ReviewProposal}, \srsref{FR_ApprovalResponse}\\
        \bottomrule
    \end{longtable}

    Some requests from interviews did not become requirements, and the SRS lists them as out of scope.
    Table~\ref{tab:excluded} gives each with the reason.

    \begin{table}[h!]
        \centering
        \renewcommand{\arraystretch}{1.2}
        \begin{tabularx}{\textwidth}{>{\raggedright\arraybackslash}p{0.33\textwidth} >{\raggedright\arraybackslash}X}
            \toprule
            \textbf{Request} & \textbf{Decision}\\
            \midrule
            Booking meetings; suggested meeting times (\intref{CI}{1}, \intref{ES}{1}, \intref{S}{1}) & Not supported. Meetings take place outside the platform.\\
            Messaging between users (\intref{ES}{1}) & Not supported. Discussion is limited to threads on posts and applications.\\
            Supervisor feedback and reports during the project (\intref{IS}{1}, \intref{SEI}{1}, \intref{CI}{1}) & Not supported. The platform covers the process up to team project approval.\\
            Resolving IP and confidentiality questions (\intref{SEI}{1}) & Not supported. This stays with the instructor and university.\\
            Matching supervisors to student projects (\intref{ES}{1}, \intref{IS}{1}) & Not supported. Student-proposed projects are visible to supervisors, and teams and supervisors find each other.\\
            Automatic approval decisions (\intref{SEI}{1}) & Never. Decisions stay with instructors and TAs.\\
            \bottomrule
        \end{tabularx}
        \caption{Requests from the interviews that were not included}
        \label{tab:excluded}
    \end{table}

    \section{Results and Next Steps}
    \label{sec:results}

    The interviews shaped the SRS in these ways.

    \begin{itemize}
        \item The SE instructor's interview set the main problems and priorities, including repetitive review,
        unselected projects, and record drift. It also led to modularity and to the principle that
        approval decisions are never automated.
        \item The TRON and CS interviews showed that courses approve projects differently. This led to the three
        approval modes (in-platform, external, integrated), configurable course rules, and the
        requirement to import enrollment from Avenue or a file.
        \item The CS instructor's data question led to the SRS rationale on hosting and to the scoped access
        for external supervisors, who are unverified until an instructor approves one of their proposals.
        \item The supervisor interviews led to the proposal template, early course information, student and
        team profiles, the recorded agreement between a supervisor and a team, and notifications with a checklist.
    \end{itemize}

    The following are still open or planned.

    \begin{itemize}
        \item \textbf{Validation of the SRS.} We will present the SRS to Dr.~Smith in a walkthrough and
        checklist session, and demonstrate the platform to Dr.~Moradi and Dr.~Wassyng. This feedback will be
        added to this report as it happens.
        \item \textbf{More interviews.} We will interview stakeholders again after the first working version,
        to check that it addresses what we heard. We should also interview TAs, and CS and TRON students,
        and ask the CS instructor to confirm team sizes and program rules.
        \item \textbf{Open questions.} How involved the platform should be in the TRON course (\intref{TI}{1}),
        and whether external supervisors will accept the unverified and verified model, remain open.
    \end{itemize}

    \section{Limitations}
    \label{sec:limits}

    \begin{itemize}
        \item \textbf{Small sample.} We interviewed one student, one internal supervisor, one joint session with two
        external supervisors, and one instructor per course. The broad claims in this report should be validated
        with more stakeholders.
        \item \textbf{Student coverage.} The one student went through the SE (GitLab) process. CS and TRON students
        are not represented. Our own team's experience as students also shaped the project, and may bias
        how we read other students' statements.
        \item \textbf{No TA interview.} TAs do much of the day-to-day work, and we have no direct evidence
        from them.
        \item \textbf{External supervisors.} The two external supervisors were first-time supervisors from a
        single company, and were interviewed together, so one may have influenced the other. Experienced sponsors
        and sponsors from other industries may see the process differently.
        \item \textbf{Notes, not transcripts.} Interviews were not recorded, and the notes record our
        understanding of what was said. We do not have original wording, and some nuance may be lost. The
        first interviews (October 1 and 7) took place before the guide was finalized and were less structured.
        \item \textbf{Interpretation.} Thinks and Feels entries in the empathy maps are our inference, and
        we built the maps after the interviews, from the notes. Interviewees have not reviewed them.
        \item \textbf{Informal process.} The interviews were part of a course project, and we did not seek
        formal ethics clearance or written consent. We withheld the names of students and supervisors, but
        instructors can be identified.
        \item \textbf{Stakeholder bias.} The SE instructor is also our supervisor, and the SE process is
        the one we know best, so the requirements may lean toward SE.
    \end{itemize}

    \newpage
    \appendix

    \section{Generative AI Use Disclosure}
    \label{app:ai-use}

    \subsection{AI Tools Used}

    \begin{itemize}
        \item \textbf{Anthropic's Claude Sonnet 5.5}
    \end{itemize}

    \subsection{How and Where Used}

    \begin{itemize}
        \item Helping to map the themes to SRS requirements and drafting the coverage table (Table~\ref{tab:trace} and
        Table~\ref{tab:coverage}).
        \item Formatting in \LaTeX{} and the BibTeX entries, especially with the table syntax.
    \end{itemize}

    \subsection{Verification of AI-Generated Content}

    \begin{itemize}
        \item The team checked the requirement names in Table~\ref{tab:trace} against the published SRS, and the
        references against their sources.
        \item The team is responsible for the accuracy, completeness, and consistency of this document.
    \end{itemize}

    \bibliographystyle {plainnat}
    \bibliography {../../../refs/References}

\end{document}